% GENERATED FILE. Edit canonical lessons, then run npm run generate:books.
% canonical-source-hash: fnv1a64:47338dc53a5515a8
% canonical-lessons: GU-C378-had, GU-C378-maryada, GU-C378-atisay, GU-C378-praman, GU-C378-khamvu, GU-R378-review-371-378

\chapter{Within Limits}
\label{ch:gu-within-limits}
\hlchaptermodality{378}

\begin{chapteropening}
\textbf{By the end of this chapter:} \emph{I can say how far is too far.}

Complete the last lesson of chapter 378: I can say how far is too far.
\end{chapteropening}

\section[\texorpdfstring{\gu{હદ} (had)}{હદ (had)}]{\gu{હદ} (had) — a limit}

\label{lesson:GU-C378-had}



\begin{quote}
Before the new one: say the Gujarati for to become different, then the Gujarati for to worsen.
\end{quote}

\subsection*{You'll want to know: \gu{હદ}}
\textbf{\gu{હદ}} — \emph{had} — \textquotedblleft{}a limit\textquotedblright{}.

\begin{quote}
\gu{જો} \gu{તમે} \gu{હદ} \gu{ઓળંગશો} \gu{તો} \gu{હું} \gu{ગુરુને} \gu{ફરિયાદ} \gu{કરીશ}.
\end{quote}

\begin{quote}
\emph{jo tame had oḷangśo to hũ gurune phariyād karīś.}
\end{quote}

\begin{quote}
If you cross the line, I will complain to the teacher.
\end{quote}

A feminine noun.

\begin{cousinweb}
\textbf{\gu{હદ}} (\emph{had}) is usually traced to Arabic \emph{ḥadd}, \textquotedblleft{}a boundary, a limit\textquotedblright{}, through Persian.
\end{cousinweb}

\begin{grammarlens}[title={What you've built}]
One more word for saying what happens if.
\end{grammarlens}

\subsection*{Your turn}
\begin{itemize}
\raggedright
  \item \emph{Say it:} \emph{had}
  \item \emph{Say it:} \emph{had}, once more
  \item \emph{Say it:} \emph{jo tame had oḷangśo to hũ gurune phariyād karīś.} — If you cross the line, I will complain to the teacher.
  \item \emph{Recall:} say the Gujarati for to become different, then the Gujarati for to worsen, then say \emph{had} again
\end{itemize}

\begin{tcolorbox}[breakable,colback=teal!4,colframe=teal!35!black,title={Before you move on}]
What does \gu{હદ} mean? (\textquotedblleft{}A limit\textquotedblright{}.) Say it once more.
\end{tcolorbox}
\section[\texorpdfstring{\gu{મર્યાદા} (maryādā)}{મર્યાદા (maryādā)}]{\gu{મર્યાદા} (maryādā) — propriety; a bound one keeps to}

\label{lesson:GU-C378-maryada}



\begin{quote}
Before the new one: say the Gujarati for to worsen, then the Gujarati for a limit.
\end{quote}

\subsection*{You'll want to know: \gu{મર્યાદા}}
\textbf{\gu{મર્યાદા}} — \emph{maryādā} — \textquotedblleft{}propriety; a bound one keeps to\textquotedblright{}.

\begin{quote}
\gu{જો} \gu{તમે} \gu{મર્યાદામાં} \gu{રહીને} \gu{બોલશો} \gu{તો} \gu{બધા} \gu{સાંભળશે}.
\end{quote}

\begin{quote}
\emph{jo tame maryādāmā̃ rahīne bolśo to badhā sā̃bhaḷśe.}
\end{quote}

\begin{quote}
If you speak within bounds, everyone will listen.
\end{quote}

A feminine noun.

\begin{cousinweb}
\textbf{\gu{મર્યાદા}} (\emph{maryādā}) is a tatsama, borrowed from Sanskrit \emph{maryādā}, \textquotedblleft{}a frontier, a bound\textquotedblright{} (Monier-Williams); Gujarati also uses it for the bounds of decent behaviour.
\end{cousinweb}

\begin{grammarlens}[title={What you've built}]
One more word for saying what happens if.
\end{grammarlens}

\subsection*{Your turn}
\begin{itemize}
\raggedright
  \item \emph{Say it:} \emph{maryādā}
  \item \emph{Say it:} \emph{maryādā}, once more
  \item \emph{Say it:} \emph{jo tame maryādāmā̃ rahīne bolśo to badhā sā̃bhaḷśe.} — If you speak within bounds, everyone will listen.
  \item \emph{Recall:} say the Gujarati for to worsen, then the Gujarati for a limit, then say \emph{maryādā} again
\end{itemize}

\begin{tcolorbox}[breakable,colback=teal!4,colframe=teal!35!black,title={Before you move on}]
What does \gu{મર્યાદા} mean? (\textquotedblleft{}Propriety; a bound one keeps to\textquotedblright{}.) Say it once more.
\end{tcolorbox}
\section[\texorpdfstring{\gu{અતિશય} (atiśay)}{અતિશય (atiśay)}]{\gu{અતિશય} (atiśay) — excessive}

\label{lesson:GU-C378-atisay}



\begin{quote}
Before the new one: say the Gujarati for a limit, then the Gujarati for propriety.
\end{quote}

\subsection*{You'll want to know: \gu{અતિશય}}
\textbf{\gu{અતિશય}} — \emph{atiśay} — \textquotedblleft{}excessive\textquotedblright{}.

\begin{quote}
\gu{જો} \gu{તમે} \gu{અતિશય} \gu{મીઠાઈ} \gu{ખાશો} \gu{તો} \gu{દાંત} \gu{બગડશે}.
\end{quote}

\begin{quote}
\emph{jo tame atiśay mīṭhāī khāśo to dā̃t bagaḍśe.}
\end{quote}

\begin{quote}
If you eat too many sweets, your teeth will go bad.
\end{quote}

\begin{cousinweb}
\textbf{\gu{અતિશય}} (\emph{atiśay}) is a tatsama, borrowed from Sanskrit \emph{atiśaya}, \textquotedblleft{}surpassing, excess\textquotedblright{} (Monier-Williams).
\end{cousinweb}

\begin{grammarlens}[title={What you've built}]
One more word for saying what happens if.
\end{grammarlens}

\subsection*{Your turn}
\begin{itemize}
\raggedright
  \item \emph{Say it:} \emph{atiśay}
  \item \emph{Say it:} \emph{atiśay}, once more
  \item \emph{Say it:} \emph{jo tame atiśay mīṭhāī khāśo to dā̃t bagaḍśe.} — If you eat too many sweets, your teeth will go bad.
  \item \emph{Recall:} say the Gujarati for a limit, then the Gujarati for propriety, then say \emph{atiśay} again
\end{itemize}

\begin{tcolorbox}[breakable,colback=teal!4,colframe=teal!35!black,title={Before you move on}]
What does \gu{અતિશય} mean? (\textquotedblleft{}Excessive\textquotedblright{}.) Say it once more.
\end{tcolorbox}
\section[\texorpdfstring{\gu{પ્રમાણ} (pramāṇ)}{પ્રમાણ (pramāṇ)}]{\gu{પ્રમાણ} (pramāṇ) — a proportion}

\label{lesson:GU-C378-praman}



\begin{quote}
Before the new one: say the Gujarati for propriety, then the Gujarati for excessive.
\end{quote}

\subsection*{You'll want to know: \gu{પ્રમાણ}}
\textbf{\gu{પ્રમાણ}} — \emph{pramāṇ} — \textquotedblleft{}a proportion\textquotedblright{}.

\begin{quote}
\gu{જો} \gu{દાળમાં} \gu{મીઠાનું} \gu{પ્રમાણ} \gu{વધારે} \gu{હશે} \gu{તો} \gu{તે} \gu{ખારી} \gu{લાગશે}.
\end{quote}

\begin{quote}
\emph{jo dāḷmā̃ mīṭhānũ pramāṇ vadhāre haśe to te khārī lāgśe.}
\end{quote}

\begin{quote}
If there is too much salt in the dal, it will taste salty.
\end{quote}

A neuter noun.

\begin{cousinweb}
\textbf{\gu{પ્રમાણ}} (\emph{pramāṇ}) is a tatsama, borrowed from Sanskrit \emph{pramāṇa}, \textquotedblleft{}a measure, a scale, a standard\textquotedblright{} (Monier-Williams).
\end{cousinweb}

\begin{grammarlens}[title={What you've built}]
One more word for saying what happens if.
\end{grammarlens}

\subsection*{Your turn}
\begin{itemize}
\raggedright
  \item \emph{Say it:} \emph{pramāṇ}
  \item \emph{Say it:} \emph{pramāṇ}, once more
  \item \emph{Say it:} \emph{jo dāḷmā̃ mīṭhānũ pramāṇ vadhāre haśe to te khārī lāgśe.} — If there is too much salt in the dal, it will taste salty.
  \item \emph{Recall:} say the Gujarati for propriety, then the Gujarati for excessive, then say \emph{pramāṇ} again
\end{itemize}

\begin{tcolorbox}[breakable,colback=teal!4,colframe=teal!35!black,title={Before you move on}]
What does \gu{પ્રમાણ} mean? (\textquotedblleft{}A proportion\textquotedblright{}.) Say it once more.
\end{tcolorbox}
\section[\texorpdfstring{\gu{ખમવું} (khamvũ)}{ખમવું (khamvũ)}]{\gu{ખમવું} (khamvũ) — to tolerate}

\label{lesson:GU-C378-khamvu}



\begin{quote}
Before the new one: say the Gujarati for excessive, then the Gujarati for a proportion.
\end{quote}

\subsection*{You'll want to know: \gu{ખમવું}}
\textbf{\gu{ખમવું}} — \emph{khamvũ} — \textquotedblleft{}to tolerate\textquotedblright{}.

\begin{quote}
\gu{જો} \gu{તે} \gu{ફરી} \gu{ખોટું} \gu{બોલશે} \gu{તો} \gu{હું} \gu{નહીં} \gu{ખમું}.
\end{quote}

\begin{quote}
\emph{jo te pharī khoṭũ bolśe to hũ nahī̃ khamũ.}
\end{quote}

\begin{quote}
If he lies again, I will not stand for it.
\end{quote}

Past, with \gu{તેણે} (he or she …): \gu{તેણે} \gu{ગુરુનો} \gu{ઠપકો} \gu{ખમ્યો}, \emph{teṇe guruno ṭhapko khamyo}.

\begin{cousinweb}
\textbf{\gu{ખમવું}} (\emph{khamvũ}) is inherited from Sanskrit \emph{kṣámatē}, \textquotedblleft{}is patient, endures\textquotedblright{}; CDIAL lists the Gujarati verb (CDIAL 3657).
\end{cousinweb}

\begin{grammarlens}[title={What you've built}]
One more word for saying what happens if.
\end{grammarlens}

\subsection*{Your turn}
\begin{itemize}
\raggedright
  \item \emph{Say it:} \emph{khamvũ}
  \item \emph{Say it:} \emph{khamvũ}, once more
  \item \emph{Say it:} \emph{jo te pharī khoṭũ bolśe to hũ nahī̃ khamũ.} — If he lies again, I will not stand for it.
  \item \emph{Recall:} say the Gujarati for excessive, then the Gujarati for a proportion, then say \emph{khamvũ} again
\end{itemize}

\begin{tcolorbox}[breakable,colback=teal!4,colframe=teal!35!black,title={Before you move on}]
What does \gu{ખમવું} mean? (\textquotedblleft{}To tolerate\textquotedblright{}.) Say it once more.
\end{tcolorbox}
\section[(dialogue)]{Looking back: from Gains and Losses to Within Limits}

\label{lesson:GU-R378-review-371-378}



\begin{quote}
Say the Gujarati for a benefit and an advantage.
\end{quote}

\subsection*{Your turn}
Say these aloud:

\begin{itemize}
\raggedright
  \item a benefit, an advantage, use, a debt, to lose for good
  \item insurance, a guarantee, a protective cover, alert, to include
  \item a way out, a solution, an alternative, a difficulty, to get in the way
  \item an obstruction, a hindrance, trouble, to keep in check, to avert
  \item essential, unavoidable, compulsory, to get done, to provide
  \item a mistake, negligence, a defect, slackness, to get better
  \item a plight, steady, normal, to become different, to worsen
  \item a limit, propriety, excessive, a proportion, to tolerate
\end{itemize}

\begin{tcolorbox}[breakable,colback=teal!4,colframe=teal!35!black,title={Before you move on}]
A benefit? (\textbf{\gu{ફાયદો}}, \emph{phāydo}.) An advantage? (\textbf{\gu{લાભ}}, \emph{lābh}.) Use? (\textbf{\gu{ઉપયોગ}}, \emph{upyog}.) A debt? (\textbf{\gu{કરજ}}, \emph{karaj}.) To lose for good? (\textbf{\gu{ખોવું}}, \emph{khovũ}.) Insurance? (\textbf{\gu{વીમો}}, \emph{vīmo}.) A guarantee? (\textbf{\gu{ગેરંટી}}, \emph{geranṭī}.) A protective cover? (\textbf{\gu{કવચ}}, \emph{kavac}.) Alert? (\textbf{\gu{સચેત}}, \emph{sacet}.) To include? (\textbf{\gu{આવરવું}}, \emph{āvarvũ}.) A way out? (\textbf{\gu{ઉપાય}}, \emph{upāy}.) A solution? (\textbf{\gu{નિરાકરણ}}, \emph{nirākaraṇ}.) An alternative? (\textbf{\gu{વિકલ્પ}}, \emph{vikalp}.) A difficulty? (\textbf{\gu{મુશ્કેલી}}, \emph{muśkelī}.) To get in the way? (\textbf{\gu{નડવું}}, \emph{naḍvũ}.) An obstruction? (\textbf{\gu{અવરોધ}}, \emph{avrodh}.) A hindrance? (\textbf{\gu{વિઘ્ન}}, \emph{vighna}.) Trouble? (\textbf{\gu{મુસીબત}}, \emph{musībat}.) To keep in check? (\textbf{\gu{ખાળવું}}, \emph{khāḷvũ}.) To avert? (\textbf{\gu{નિવારવું}}, \emph{nivārvũ}.) Essential? (\textbf{\gu{આવશ્યક}}, \emph{āvaśyak}.) Unavoidable? (\textbf{\gu{અનિવાર્ય}}, \emph{anivārya}.) Compulsory? (\textbf{\gu{ફરજિયાત}}, \emph{pharjiyāt}.) To get done? (\textbf{\gu{પતવું}}, \emph{patvũ}.) To provide? (\textbf{\gu{પૂરું} \gu{પાડવું}}, \emph{pūrũ pāḍvũ}.) A mistake? (\textbf{\gu{ભૂલ}}, \emph{bhūl}.) Negligence? (\textbf{\gu{ગફલત}}, \emph{gaphlat}.) A defect? (\textbf{\gu{ખામી}}, \emph{khāmī}.) Slackness? (\textbf{\gu{ઢીલ}}, \emph{ḍhīl}.) To get better? (\textbf{\gu{સુધરવું}}, \emph{sudharvũ}.) A plight? (\textbf{\gu{દશા}}, \emph{daśā}.) Steady? (\textbf{\gu{સ્થિર}}, \emph{sthir}.) Normal? (\textbf{\gu{સામાન્ય}}, \emph{sāmānya}.) To become different? (\textbf{\gu{બદલાવું}}, \emph{badlāvũ}.) To worsen? (\textbf{\gu{વણસવું}}, \emph{vaṇasvũ}.) A limit? (\textbf{\gu{હદ}}, \emph{had}.) Propriety? (\textbf{\gu{મર્યાદા}}, \emph{maryādā}.) Excessive? (\textbf{\gu{અતિશય}}, \emph{atiśay}.) A proportion? (\textbf{\gu{પ્રમાણ}}, \emph{pramāṇ}.) To tolerate? (\textbf{\gu{ખમવું}}, \emph{khamvũ}.)
\end{tcolorbox}
